\section{与已有理论的关系}
\label{sec:comparison}

本节先说明指数 $\gamma=1$ 时 RL 条件与若干二次型条件的关系，
再比较输出误差界给出的距离收缩上界，最后用第~
\ref{sec:rank-deficient-hamilton}~节的 Hamilton 实例说明
完整邻域上若干指定证书的失效。所引算子框架均允许集值映射；
本文的区别在于所检验不等式的阶数及其与输出误差界的配合。

\subsection{指数为一时的二次型条件}

\begin{proposition}[指数为一时的条件识别]
\label{prop:linear-identifications}
固定图像块 $\mathcal G$、尺度 $R$ 和 $\lambda>0$。
对任意 $(u,u^*),(v,v^*)\in\mathcal G$，记
$a=u-v$、$b=u^*-v^*$。
以下结论均限于满足 $\|a+\lambda b\|\leq R$ 的点对。
\begin{enumerate}
\item $\mathrm{RL}(1,L)$ 等价于
\begin{equation}
 \lambda\langle a,b\rangle
 \geq \frac{1-L^2}{2(1+L^2)}
       \bigl(\|a\|^2+\lambda^2\|b\|^2\bigr).
 \label{eq:linear-rl-balanced}
\end{equation}
因此，$\mathrm{RL}(1,1)$ 恰为这些点对上的单调性。
采用半单调性不等式
$\langle a,b\rangle\geq\mu\|a\|^2+\nu\|b\|^2$，
其参数为
\begin{equation}
 \mu=\frac{1-L^2}{2\lambda(1+L^2)},
 \qquad \nu=\lambda^2\mu.
 \label{eq:balanced-slice-parameters}
\end{equation}
故 $\mathrm{RL}(1,L)$ 对应该不等式的平衡参数切片。

\item 对任意 $\tau\geq0$，尺度化次单调性条件
\[
 \lambda\langle a,b\rangle
 \geq-\tau\|a+\lambda b\|^2
\]
等价于 $\mathrm{RL}(1,L)$，其中
\begin{equation}
 L^2=1+4\tau.
 \label{eq:luke-tam-normalization}
\end{equation}

\end{enumerate}
\end{proposition}

\begin{proof}
当 $\gamma=1$ 时，RL 内积不等式为
\[
 4\lambda\langle a,b\rangle
 \geq(1-L^2)\|a+\lambda b\|^2.
\]
展开平方并合并内积项，即得
\eqref{eq:linear-rl-balanced}。
令 $L=1$ 得到单调性；两边除以 $\lambda$ 得到
\eqref{eq:balanced-slice-parameters}。
代入 $1-L^2=-4\tau$，即得第二项。

\end{proof}

式~\eqref{eq:balanced-slice-parameters} 使用
\cite{EvensPasLatafatPatrinos2025SemimonotonePPA}
的半单调性约定。
式~\eqref{eq:luke-tam-normalization} 则识别了
Luke--Tam 的尺度化次单调性条件
\cite[Definition~2]{LukeTam2025GeneralizedMonotonicity}。
当 $\tau<1/2$ 时，该文命题~4 将截断图像上的条件等价地写成
相应反射预解式的 Lipschitz 估计，常数为 $\sqrt{1+4\tau}$。
该文假设~2(d)(ii) 的局部 Minty 值域覆盖可以供应本文 (B1)，
具体映射见备注~\ref{rem:LT-coverage-map}。
本文 (B2) 允许非线性阶数及随输出选择的最近零锚，
(B3)--(B4) 则把输出真残差误差界与该阶数相配合。

\subsection{一般半单调性与锚定条件的接口}

平衡切片以外的半单调性也能供应线性 RL。
仍固定同一图块及受检点对，设
\begin{equation}
 \langle a,b\rangle\geq\mu\|a\|^2+\nu\|b\|^2,
 \qquad \mu,\nu\in\mathbb R.
 \label{eq:general-semimonotonicity}
\end{equation}
令
\[
 A=1+\lambda\mu+\nu/\lambda,\qquad
 B=\lambda\mu-\nu/\lambda,\qquad
 \Delta=1-4\mu\nu.
\]
当 $A>0$、$\Delta\geq0$ 时，该条件给出
\begin{equation}
 \mathrm{RL}(1,L_C),\qquad
 L_C=\frac{|B|+\sqrt\Delta}{A}.
 \label{eq:semimonotone-cayley-factor}
\end{equation}
这个转接只需配方。写 $d=a+\lambda b$、$c=a-\lambda b$，
式~\eqref{eq:general-semimonotonicity} 等价于
\[
 A\|c\|^2+2B\langle d,c\rangle
 \leq(1-\lambda\mu-\nu/\lambda)\|d\|^2,
\]
亦即 $\|Ac+Bd\|^2\leq\Delta\|d\|^2$。
三角不等式随即给出 \eqref{eq:semimonotone-cayley-factor}。
特别地，以下条件均在相同点对范围内供应线性 RL：
\begin{center}
\begin{tabular}{lcc}
条件 & 步长门 & 反射常数 \\
\hline
单调性 & $\lambda>0$ & $1$ \\
$h$-亚单调性 & $h\geq0,\ \lambda h<1$ &
 $(1+\lambda h)/(1-\lambda h)$ \\
$\eta$-负余单调性 & $\eta\geq0,\ \lambda>\eta$ &
 $(\lambda+\eta)/(\lambda-\eta)$
\end{tabular}
\end{center}
半单调性及其预解式参数门已有系统研究
\cite[Definition~4.1, Proposition~4.12]{EvensEtAl2024SemimonotoneV2}；
这里使用该 arXiv v2 的约定和编号，不将转接公式作为新颖性主张。
全图同参数极大性在文献的可行步长域内可进一步保证满输入域；
局部图块的配方本身仅给唯一性，不给输入覆盖。
边界 $A=0$ 不能代入公式：例如 $F(y)=-y/\lambda$、
$\mu=-1/(2\lambda)$、$\nu=-\lambda/2$ 满足上述不等式，
但其 Minty 输入域只有 $\{0\}$，该输入的近端输出为整个空间。

同一配方也适用于一个零锚 $(p,0)$。
若每个允许输出都能以输入的某个最近零点为锚满足
式~\eqref{eq:general-semimonotonicity}，它便供应 (B2)，
而无需输出之间的全对比较。
标量 weak Minty 条件
$\langle y-p,v\rangle\geq-\eta\|v\|^2$
是其中 $(\mu,\nu)=(0,-\eta)$ 的情形；
文献中的固定矩阵、固定零锚集合和预条件更新还带其他量词与算法门
\cite[Definition~2.1, Theorem~2.4]{EvensEtAl2024SemimonotoneV2}。
仅有某个固定零点的锚定证书，不保证该点是输入最近零点，
也不自动提供全对唯一性、(B1)、(B3)--(B4) 或留域预算。

\subsection{输出误差界与距离收缩上界}

本文的正则性条件作用于近端输出处的原算子：
\[
 d(y,S)\leq\psi\bigl(r_F(y)\bigr).
\]
这一条件可独立于 RL 验证。例如，若在所需输出区域上已有
度量 $q$-次正则性估计
\begin{equation}
 d(y,S)\leq K r_F(y)^q,
 \qquad K>0,\quad q>0,
 \label{eq:power-eb-comparison}
\end{equation}
则可取 $\psi(t)=Kt^q$。调用已有误差界时须保持同一目标集和输出区域：
若标准次正则性只给距离到 $F^{-1}(0)$，而本文选择了严格子集
$S\subset F^{-1}(0)$，则不能自动得到关于 $S$ 的估计。
Li--Mordukhovich 的定义允许 $q>0$，其最大单调 PPA 应用使用
$0<q\leq1$ \cite[Definition~3.1, Section~7]{LiMordukhovich2012HolderSubregularityPPA}。
Mordukhovich--Ouyang 进一步研究 $q>1$ 的高阶次正则性及
Newton 类应用 \cite{MordukhovichOuyang2015HigherOrderSubregularity}；
这些正则性结果供应的是 (B3)，不能替代当前 PPA 的其他前件。
当 $0<\gamma<1$、$L>0$ 时，直接递推上界的右端满足
\[
 K\left(\frac{d+Ld^\gamma}{2\lambda}\right)^q
 =K\left(\frac{L}{2\lambda}\right)^q
   d^{\gamma q}(1+o(1))\qquad(d\downarrow0).
\]
因此，$\gamma q>1$ 时该上界在充分小尺度上满足 (B4)；
当 $\gamma q=1$ 时，还须
$K(L/(2\lambda))^q<1$。
当 $\gamma q<1$ 或临界系数不小于一时，该直接上界不能给出
统一严格收缩。这是当前证书的边界，不是算法发散的判据，
也不能据此认定轨道的精确速率。

一般控制函数也用于 arXiv 初版
\cite[Theorems~4.3--4.4]{LiMordukhovichZhu2024Preprint}
的抽象收敛框架。该框架针对带目标函数的次梯度序列，
使用目标下降、步长替代量、相对次梯度误差、连续性及聚点等条件，
并将作用于距离的控制函数与其他模量组合成速率映射。
本文以 RL 提供步长估计，以近端包含关系提供输出残差，
再由初始半径预算证明局部化和极限存在；不要求预先给定
目标下降结构或聚点。两者在适当的次梯度 PPA 特例中可以交叠，
但现有假设本身不给出一般的相互包含关系。
第~\ref{sec:rank-deficient-hamilton}~节的误差界
$\psi(t)=3t^2$ 则由该实例的残差直接计算得到。

在 $\gamma=1$ 且输出误差界为 $\psi(t)=\rho t$、$\rho>0$
时，令 $d=d(x,S)$、$d^+=d(x^+,S)$、$s=\|x-x^+\|$。
由一步能量估计及 $d^+\leq\rho s/\lambda$，有
\[
 \left(1+\frac{\lambda^2}{\rho^2}\right)(d^+)^2
 \leq(d^+)^2+s^2
 \leq\frac{1+L^2}{2}d^2.
\]
故得到距离的一步上界
\begin{equation}
 d^+\leq\kappa_E d,
 \qquad
 \kappa_E=\sqrt{\frac{1+L^2}{2}}\,
          \frac{\rho}{\sqrt{\rho^2+\lambda^2}}.
 \label{eq:linear-output-coupling}
\end{equation}
当 $L=1$ 时，系数退化为
$\rho/\sqrt{\rho^2+\lambda^2}$，与最大单调情形下
Leventhal 定理~3.1 的距离上界具有相同形式
\cite{Leventhal2009MetricSubregularityPPA}。

这条能量上界与主定理的直接兼容门须区分。
线性误差界下，(B4) 要求
\[
 \kappa_D=\frac{\rho(1+L)}{2\lambda}<1,
\]
它不由 $\kappa_E<1$ 自动推出。
例如，取最大单调 $F(y)=y/2$、$\lambda=1$、$S=\{0\}$，
则 $J_F(x)=2x/3$，最小线性 RL 常数 $L=1/3$、最小 EB 常数 $\rho=2$。
此时 $\kappa_D=4/3>1$，任意缩小正半径都不能满足 (B4)，
而 $\kappa_E=2/3$ 恰为实际距离因子。
因此，本文直接分支未包含经典最大单调定理的全部步长范围。
以下比较仅针对能量的一步系数；本文未据此另行声明一个
省去 (B4) 的局部收敛分支。

采用 Luke--Tam 的尺度化参数 $L^2=1+4\tau$，有
\begin{equation}
 \kappa_E^2=(1+2\tau)\frac{\rho^2}{\rho^2+\lambda^2},
 \qquad
 \kappa_E<1\quad\Longleftrightarrow\quad
 2\tau\rho^2<\lambda^2.
 \label{eq:coupled-linear-factor}
\end{equation}
Luke--Tam 定理~2 在其假设~2 下给出的保证系数为
\cite{LukeTam2025GeneralizedMonotonicity}
\[
 \kappa_{\rm LT}^2
 =1+2\tau-\left(\frac{\lambda}{\lambda+\rho}\right)^2,
 \qquad
 \tau(1+\rho/\lambda)^2<\frac12.
\]
对共同的尺度化次单调性与误差界参数，保留输出端误差界的耦合
给出较小的保证系数，以及较弱的上述标量收缩限制。
这只比较相同一步不等式下的上界；局部收敛还须检验值域覆盖、
输出误差界适用区域和相应的初值局部化条件，
不能据此认定两条局部定理的全部假设存在包含关系。

\subsection{Hamilton 实例排除的指定条件}

近乎平均映射理论
\cite{LukeThaoTam2018Quantitative} 与广义单调性 PPA 理论
\cite{LukeTam2025GeneralizedMonotonicity,
EvensPasLatafatPatrinos2025SemimonotonePPA}
均允许集值算子。以下比较限于同一完整邻域上的指定证书。

\begin{proposition}[指定二次型证书在完整邻域上的失效]
\label{prop:hamilton-quadratic-separation}
令 $F$ 为第~\ref{sec:rank-deficient-hamilton}~节的算子，
$J=J_F$，并固定任意零点 $\bar z=((0,\bar s),0)\in S$。
\begin{enumerate}
\item 在 $\bar z$ 的任意完整输入邻域上，不存在有限常数
$C\geq0$，使该邻域内所有输入 $z,w$ 均满足
\begin{equation}
 \|Jz-Jw\|^2+
 \|(z-Jz)-(w-Jw)\|^2\leq C\|z-w\|^2.
 \label{eq:excluded-quadratic-J}
\end{equation}
\item 在 $\bar z$ 的任意完整输入邻域上，不存在有限常数
$C\geq0$，使所有输入 $z$ 均满足
$\|Jz-\bar z\|\leq C\|z-\bar z\|$。
\item 在 $(\bar z,0)$ 的任意完整图像邻域上，
不存在有限的亚单调性模量。
\item 不存在固定有限矩阵 $V$，使所有充分接近
$(\bar z,0)$ 的图像点 $(u,u^*)$ 均满足
\begin{equation}
 \langle u-\bar z,u^*\rangle\geq\langle u^*,Vu^*\rangle.
 \label{eq:excluded-fixed-anchor-matrix}
\end{equation}
\end{enumerate}
上述反例亦可取在命题~\ref{prop:rank-deficient-hamilton-rleb}
使用的图像块内。
\end{proposition}

\begin{proof}
对 $t>0$，取输入 $w_t=((0,\bar s),(0,t))$。
由近端方程~\eqref{eq:hd-proximal-equations}，
\[
 Jw_t=\bigl((0,\bar s+\sqrt{t/100}),(0,t/100)\bigr),
 \qquad J\bar z=\bar z.
\]
于是
\[
 \|w_t-\bar z\|^2=t^2,
 \qquad
 \|Jw_t-J\bar z\|^2+
 \|(w_t-Jw_t)-(\bar z-J\bar z)\|^2
 =\frac{t}{50}+\frac{4901}{5000}t^2.
\]
与 $t^2$ 的比值趋于 $+\infty$，证明第一项。
同时
\[
 \frac{\|Jw_t-\bar z\|}{\|w_t-\bar z\|}
 =\sqrt{\frac1{100t}+\frac1{10000}}
 \longrightarrow+\infty,
\]
证明第二项。

对第三项，取图像点
\[
 u_t=((0,\bar s+t),(0,t)),
 \qquad
 u_t^*=F(u_t)=((0,-\sqrt t),(0,99t)).
\]
将其与 $(\bar z,0)$ 比较，有
\[
 \frac{\langle u_t-\bar z,u_t^*\rangle}
      {\|u_t-\bar z\|^2}
 =-\frac1{2\sqrt t}+\frac{99}{2}
 \longrightarrow-\infty.
\]
故任何有限亚单调性模量均失效。

最后，改取 $u_t=((0,\bar s+t^{1/4}),(0,t))$，
$u_t^*$ 保持不变。则
\[
 \frac{\langle u_t-\bar z,u_t^*\rangle}{\|u_t^*\|^2}
 =\frac{-t^{3/4}+99t^2}{t+9801t^2}
 \longrightarrow-\infty.
\]
而对任何固定有限矩阵 $V$，
$\langle u_t^*,Vu_t^*\rangle/\|u_t^*\|^2\geq-\|V\|$，
故第四项成立。各点均趋于 $(\bar z,0)$。
第一、二项使用的近端输入到 $S$ 的距离为 $t$；
第三、四项图像点对应的近端输入到 $S$ 的距离为 $100t$。
因此，充分小的 $t$ 也使这些反例属于所指定的图像块。
\end{proof}

例如，第一项排除有限常数的全点对 $J$ 侧二次型估计，
因而排除同一完整图像邻域上的 Luke--Tam 次单调性条件。
该结论不比较整个半单调、弱 Minty、近乎平均或部分正则性框架；
这些框架可以采用不同的相对区域、锚点测试或退化参数。
本例在相同输入管状区域上仍具有 $1/2$ 阶 RL 条件与
二次输出误差界，两者满足 RLEB 的相容性判据。

对于 Hamilton 下降法，需进一步区分更新形式与已有收敛命题。
Maddison 等的 arXiv 初版~\cite{MaddisonEtAl2018HamiltonianDescent}
命题~3.2 针对全隐式更新，使用假设 A、B，
其中包括唯一极小点、阻尼 $\Gamma\in(0,1)$、
动能对中心共轭能量的支配，以及统一能量配对界
\[
 |\langle\nabla f(x),\nabla k(p)\rangle|
 \leq C_{f,k}\bigl(f(x)-f(x^*)+k(p)\bigr)
 \qquad(x,p\in\mathbb R^d).
\]
该命题还要求初始动量 $p^0=0$ 和步长
$0<\varepsilon<(1-\Gamma)/(2\max\{C_{f,k},1\})$，
所得结论为目标函数值的收敛上界。
本例的极小点不唯一且 $\Gamma=99$。
即使只检验能量配对界，取 $x=p=(t,0)$ 也有
\[
 \frac{|\langle\nabla f(x),\nabla k(p)\rangle|}{f(x)+k(p)}
 =\frac{3}{4\sqrt t}\longrightarrow+\infty
 \qquad(t\downarrow0),
\]
故该统一界在任何完整局部邻域上亦不成立。
因此，前节对步长为一、允许非零初始动量的整条位置--动量
轨道给出的局部收敛结论，不能直接由该命题推出。
这一比较限于 arXiv 初版的命题~3.2，不表示已有 Hamilton
下降方法的其他结果均不适用。
